%% kapa: exam=1 round=37 year=2026 period=2 form=A subject=accounting subject_ko=회계학 seq=12 number=52
%% source: Q-Net 감정평가사 기출문제 (cst00309) artlSeq=5258924 file=raw/1008/5258924/2026년 제37회 감정평가사 1차 2교시.pdf
%% conversion: kordoc 4.24.51 -> tools/gen1.py (auto); review=proofread-sonnet; flags=none
\begin{question}{52}
(주)감평은 20×1년 초 액면금액을 만기일에 일시상환하는 조건인 일반사채(액면금액 ￦100,000, 표시이자는 매년 말 지급, 만기 3년, 유효이자율법으로 이자비용 인식)를 발행하였다. 동 사채는 (주)감평의 유일한 사채이며, 사채 발행 이후 조건의 변경은 없다. 동 사채에 대한 연도별 이자 관련 정보가 다음과 같을 때, 동 사채의 발행시 유효이자율은? (단, 단수차이로 인한 오차가 있으면 가장 근사치를 선택한다.)
\par
{\small\centering
\begin{tabular}{|c|r|r|r|}\hline
구분 & \multicolumn{1}{c|}{20×1년} & \multicolumn{1}{c|}{20×2년} & \multicolumn{1}{c|}{20×3년} \\\hline
연도말 표시이자지급액 & ￦10,000 & ￦10,000 & ￦10,000 \\\hline
포괄손익계산서상 이자비용 & 13,288 & 13,781 & 14,348 \\\hline
\end{tabular}\par}
\begin{choices}
  \item 연 9 \%
  \item 연 10 \%
  \item 연 11 \%
  \item 연 13 \%
  \item 연 15 \%
\end{choices}
\end{question}
\answer{5}
